% ======================================================================
\section{Perturbative expansion of the evolution operator \texorpdfstring{$\Omega$}{Omega}}\label{sec:omega}
% ======================================================================

In this section we determine the coefficients of the expansion~\eqref{eq:OmegaExp} to the order needed for the single inclusive gluon cross section at next-to-leading order. We first collect the results for the soft vacuum as initial state, which are known~\cite{Lublinsky:2016meo}. We then compute the dressed one-gluon state in LCPT, component by component, and read off the coefficients $\mf A_1$, $\mf B_3$ and $\mf C_1$. Finally we impose the unitarity of $\Omega$, derive the relations among the coefficients order by order in $g$, and summarise the coefficients in momentum and in coordinate space.

In the diagrams of this section the straight horizontal line represents the valence charges, curly lines are soft gluons, and a vertical dashed line marks an intermediate state whose energy enters a denominator. A component of the dressed one-gluon state is denoted by $\ket{(\Psi^a_i(k))_{\dots}}$, where $(k,a,i)$ are the momentum, colour and polarisation of the incoming gluon and the subscript lists the outgoing gluons and the charges, as explained below Eq.~\eqref{eq:fock}. Outgoing and intermediate momenta are integrated over.

\subsection{The vacuum incoming state}\label{sec:vacin}

The dressing of the soft vacuum, $\ket{\Psi_0}=\Omega\ket0$, was computed to order $g^3$ in~\cite{Lublinsky:2016meo}, including quarks. We do not repeat the derivation and quote the results that are needed, restricted to pure Yang--Mills theory and written in the conventions of Section~\ref{sec:setup}. They follow from the third-order series~\eqref{eq:pt3} (Appendix~\ref{app:pt}) with the matrix elements~\eqref{eq:m1}--\eqref{eq:m5} (Appendix~\ref{app:matrixelements}); the coefficients are read off by comparison with Eq.~\eqref{eq:Omegavac}.

\paragraph{Order $g$.} The valence charge emits one soft gluon,
\begin{align}\label{eq:PsiLO}
\ket{\Psi^{\rm LO}_{g\rho}}=-\int_\Lambda^\vee dk^+\!\int d^2k\,\ket{g^a_i(k)}\frac{\braket{g^a_i(k)|H_g|0}}{E_g(k)}
=-\int_\Lambda^\vee dk^+\!\int d^2k\,\frac{g\,k^i\rho^a(-\v k)}{2\pi^{3/2}\sqrt{k^+}\,\v k^2}\ket{g^a_i(k)}\,,
\end{align}
so that, by Eq.~\eqref{eq:Omegavac},
\begin{align}\label{eq:A1mom}
\mf A^{(1)a}_{2i}(k)=-\frac{\sqrt2\,g\,k^i\rho^a(-\v k)}{\sqrt{k^+}\,\v k^2}\,,\qquad \Lambda<k^+<\vee\,.
\end{align}
This is the Weizs\"acker--Williams field of the valence charge. The coefficient of the vacuum component at order $g^2$ is the normalisation term of Eq.~\eqref{eq:pt3},
\begin{align}\label{eq:N2mom}
\mf N_2=-\frac12\braket{\Psi^{\rm LO}_{g\rho}|\Psi^{\rm LO}_{g\rho}}=-\frac{g^2}{8\pi^3}\,\log\frac{\vee}{\Lambda}\int d^2p\,\frac{\rho^a(-\v p)\,\rho^a(\v p)}{\v p^2}\,.
\end{align}
It is Hermitian, and it does not commute with a single charge.

\paragraph{Order $g^2$: two gluons.} The two-gluon component arises from the emission of two gluons by the valence charge, from the emission of one gluon that splits, and from the instantaneous emission of a pair. The product of two charges in the first contribution is decomposed with Eq.~\eqref{eq:decomp}. The resulting coefficient $\mf B_2$~\cite{Lublinsky:2016meo} is given in Eq.~\eqref{eq:B2mom} of Appendix~\ref{app:coefficients}. It is the sum of a part linear in $\rho$ (commutator of the two emissions, splitting, instantaneous emission), with support $k^++p^+<\vee$ because the parent gluon is a soft mode, and of the anticommutator part $\propto\{\rho^b(-\v p),\rho^a(-\v k)\}$, with support $\Lambda<k^+,p^+<\vee$: two independently emitted gluons are each soft modes, with no condition on $k^++p^+$.

\NOTE{The commutator part of the two-emission diagram is written with the support $k^++p^+<\vee$; by its kinematics it could also be the full square $\Lambda<k^+,p^+<\vee$. To be decided by the authors.}

\paragraph{Order $g^3$: one gluon.} The one-gluon component at order $g^3$ defines $\mf A^{(3)}_2$. It contains terms with one, two and three charges,
\begin{align}
\mf A^{(3)}_2=\big[\mf A^{(3)}_2\big]_\rho+\big[\mf A^{(3)}_2\big]_{\rho\rho}+\big[\mf A^{(3)}_2\big]_{\rho\rho\rho}\,,
\end{align}
which originate from the diagrams with one, two and three valence vertices and from the two normalisation terms of Eq.~\eqref{eq:pt3}, after the decomposition~\eqref{eq:decomp}. The transverse loop integral of the self-energy of the emitted gluon is ultraviolet divergent; it is done in $d=2-\epsilon$ transverse dimensions, and $\mu_{\MS}$ is the $\MS$ scale. In the limit $\Lambda\to0$ wherever it can be taken, the part linear in $\rho$ is
\begin{align}\label{eq:A3onerho}
\Big[\mf A^{(3)a}_{2i}(k)\Big]_{\rho}=\ &\frac{g^3N_c\,\rho^a(-\v k)\,k^i}{8\sqrt{2k^+}\,\pi^2\,\v k^2}\Bigg\{\Big[-\frac2\epsilon+\log\frac{\v k^2}{\mu^2_{\MS}}\Big]
\Big[\frac{11}{3}-2\log\frac{k^+}{\Lambda}+\log\Big(1-\frac{k^+}{\vee}\Big)\Big]
\notag\\&\hspace{1.5cm}
+\frac{\pi^2}{3}-\frac{67}{9}+\log^2\Big(1-\frac{k^+}{\vee}\Big)+\log\Big(1-\frac{k^+}{\vee}\Big)\log\frac{\vee}{k^+}\Bigg\}\,.
\end{align}
The part with two charges is proportional to $f^{abc}\{\rho^b(-\v p),\rho^c(-\v k+\v p)\}$ and is given in Eq.~\eqref{eq:A3tworho} of Appendix~\ref{app:coefficients}; the part with three charges is
\begin{align}\label{eq:A3threerho}
\Big[\mf A^{(3)a}_{2i}(k)\Big]_{\rho\rho\rho}=\log\frac{\vee}{\Lambda}\,\frac{\sqrt2\,g^3\,k^i}{8\pi^3\sqrt{k^+}\,\v k^2}\,\rho^a(-\v k)\int d^2p\,\frac{\rho^b(\v p)\,\rho^b(-\v p)}{\v p^2}
=\mf A^{(1)a}_{2i}(k)\,\mf N_2\,.
\end{align}
The terms proportional to $1/\Lambda$ and $1/\vee$ of the individual diagrams cancel in Eq.~\eqref{eq:A3threerho}. The order of the charges in Eq.~\eqref{eq:A3threerho}, with $\rho^a(-\v k)$ to the left of the norm factor, follows from the order of the factors in the last term of Eq.~\eqref{eq:pt3}. It matters, because $\mf N_2$ and $\mf A^{(1)}_2$ do not commute; see Section~\ref{sec:unitarity}.

Eqs.~\eqref{eq:A3onerho}--\eqref{eq:A3threerho} correspond to Eqs.~(3.97)--(3.103) of~\cite{Lublinsky:2016meo} (arXiv version~2), with one difference: in Eq.~(3.102) there the three charges are written in the order $\rho^b\rho^b\rho^a$, whereas the order that follows from Eq.~\eqref{eq:pt3} is that of Eq.~\eqref{eq:A3threerho}.

\subsection{Calculation of the single-gluon incoming state}\label{sec:onegin}

We now compute the state $\ket{\Psi^a_i(k)}=\Omega\ket{g^a_i(k)}$ that develops from a soft gluon of momentum $k=(k^+,\v k)$, colour $a$ and polarisation $i$, with $\Lambda<k^+<\vee$. We use Eq.~\eqref{eq:pt2} with $\ket i=\ket{g^a_i(k)}$ and $E_i=E_g(k)$. The possible final states are the vacuum, one gluon, two gluons and three gluons. For each of them we give the LCPT expression, evaluate it with the matrix elements of Appendix~\ref{app:matrixelements}, and give the result in momentum space.

The long intermediate expressions, the Fourier transforms and the forms of all components in transverse coordinate space are collected in Appendix~\ref{app:onegluon}; the cancellation of the double pole, the algebra of the terms with two valence vertices and the normalisation are in Appendices~\ref{app:doublepole}, \ref{app:tworho} and~\ref{app:norm}.

\subsubsection{Vacuum outgoing state}\label{sec:vacout}

At order $g$ the incoming gluon can be absorbed by the valence charge (Figure~\ref{fig:g_vac}). This is the first-order term of Eq.~\eqref{eq:pt2} with $\ket n=\ket0$ and the denominator $E_g(k)-E_0=E_g(k)$,
\begin{align}\label{eq:psi_rho}
\ket{(\Psi^a_i(k))_\rho}=\ket0\,\frac{\braket{0|H_g|g^a_i(k)}}{E_g(k)}=\frac{g\,k^i\rho^a(\v k)}{2\pi^{3/2}\sqrt{k^+}\,\v k^2}\ket0\,,
\end{align}
where the Hermitian conjugate of Eq.~\eqref{eq:m1} was used. Comparing with the first term of Eq.~\eqref{eq:Omegag} gives
\begin{align}\label{eq:A11mom}
\mf A^{(1)a}_{1i}(k)=\frac{\sqrt2\,g\,k^i\rho^a(\v k)}{\sqrt{k^+}\,\v k^2}=-\big[\mf A^{(1)a}_{2i}(k)\big]^\dagger\,,
\end{align}
where the last equality refers to Eq.~\eqref{eq:A1mom}. The absorption amplitude is minus the conjugate of the emission amplitude, as required by unitarity (Section~\ref{sec:unitarity}). The vacuum component of order $g^3$, $\mf A^{(3)}_1$, will not be computed diagrammatically; it follows from unitarity, Eq.~\eqref{eq:A31unit}.

\begin{figure}[t]
\centering
\tikzfig{g_tovacum}
\caption{Absorption of the incoming gluon by the valence charge, Eq.~\eqref{eq:psi_rho}.}
\label{fig:g_vac}
\end{figure}

\subsubsection{One outgoing gluon}\label{sec:oneout}

At order $g^2$ the incoming gluon can turn into a gluon of different momentum $p=(p^+,\v p)$, colour $b$ and polarisation $j$ by exchanging momentum with the valence charges. The sum of the contributions defines $\mf B_3$ through the third term of Eq.~\eqref{eq:Omegag}. We distinguish the terms with one valence vertex, which are linear in $\rho$, from the terms with two valence vertices. Since the momentum exchanged with the valence modes is not zero, the outgoing gluon is never identical to the incoming one in these terms; the component along $\ket{g^a_i(k)}$ itself is the normalisation, which is discussed in Section~\ref{sec:summarywf}.

\paragraph{Single $\rho$ part: instantaneous interaction.} The instantaneous interaction~\eqref{eq:m5} contributes at first order of Eq.~\eqref{eq:pt2} (Figure~\ref{fig:g_g_one}a),
\begin{align}\label{eq:psi_inst_def}
\ket{(\Psi^a_i(k))_{g\rho\,{\rm inst}}}=\int dp^+\!\int d^2p\ \ket{g^b_j(p)}\,\frac{\braket{g^b_j(p)|H_{gg\text{-inst}}|g^a_i(k)}}{E_g(k)-E_g(p)}\,.
\end{align}
With $E_g(k)-E_g(p)=(p^+\v k^2-k^+\v p^2)/(2k^+p^+)$,
\begin{align}\label{eq:psi_inst}
\ket{(\Psi^a_i(k))_{g\rho\,{\rm inst}}}=\frac{ig^2f^{abc}}{(2\pi)^3}\int dp^+\!\int d^2p\ \frac{\sqrt{p^+k^+}\,(k^++p^+)}{(p^+-k^+)^2}\,\frac{\rho^c(\v k-\v p)}{p^+\v k^2-k^+\v p^2}\,\ket{g^b_i(p)}\,.
\end{align}
The longitudinal momentum $p^+-k^+$ is exchanged with the valence modes. The range of $p^+$ consists of two intervals (Figure~\ref{fig:kinem}),
\begin{align}\label{eq:ranges}
k^++\Lambda<p^+<\vee\quad(\text{region }1)\,,\qquad\qquad \Lambda<p^+<k^+-\Lambda\quad(\text{region }2)\,.
\end{align}
In region~1 the valence modes transfer longitudinal momentum to the gluon, in region~2 the gluon transfers longitudinal momentum to the valence modes. We denote the two contributions by $\ket{(\Psi^a_i(k))^{1}_{g\rho\,{\rm inst}}}$ and $\ket{(\Psi^a_i(k))^{2}_{g\rho\,{\rm inst}}}$; they have the same integrand. The denominator $p^+\v k^2-k^+\v p^2$ vanishes inside the region of integration and is a principal value, Eq.~\eqref{eq:PV}. \begin{figure}[t]
\centering
\tikzfig{kinem}
\caption{Longitudinal momenta allowed in the transition of a soft gluon of momentum $k^+$ into a soft gluon of momentum $p^+$ on the valence charges. One shaded region is the condition $\Lambda<|p^+-k^+|<\vee$ on the exchanged momentum, the other the conditions $\Lambda<k^+<\vee$ and $\Lambda<p^+<\vee$ on the two gluons; the allowed region is their intersection, Eq.~\eqref{eq:ranges}.}
\label{fig:kinem}
\end{figure}

\begin{figure}[t]
\centering
\begin{tabular}{ccc}
\scalebox{0.7}{\tikzfig{gtog_inst}}&\scalebox{0.7}{\tikzfig{gandgtog}}&\scalebox{0.7}{\tikzfig{gto3g_torho}}\\[2mm]
(a)&(b)&(c)
\end{tabular}
\caption{Transition of the incoming gluon into a gluon of different momentum, terms linear in $\rho$: (a) instantaneous interaction, Eq.~\eqref{eq:psi_inst_def}; (b) emission of a gluon by the valence charge followed by the merging of the two gluons, Eq.~\eqref{eq:psi1_def}; (c) splitting of the incoming gluon followed by the absorption of one of the two gluons, Eq.~\eqref{eq:psi2_def}.}
\label{fig:g_g_one}
\end{figure}

\paragraph{Single $\rho$ part: three-gluon vertex.} At second order of Eq.~\eqref{eq:pt2} there are two contributions linear in $\rho$, which contain the three-gluon vertex~\eqref{eq:m3} and one valence vertex. In the first one (Figure~\ref{fig:g_g_one}b) the valence charge emits a gluon of momentum $r=p-k$, and the two gluons merge. The intermediate state contains the gluons $k$ and $r$, so that the first denominator is $E_g(k)-E_{gg}(k,r)=-E_g(r)$:
\begin{align}\label{eq:psi1_def}
\ket{(\Psi^a_i(k))^1_{g\rho}}=\frac12\int d^3p\,d^3q\,d^3r\ \ket{g^b_j(p)}\,
\frac{\braket{g^b_j(p)|H_{ggg}|g^c_l(q)g^d_m(r)}\braket{g^c_l(q)g^d_m(r)|H_g|g^a_i(k)}}{\big(E_g(k)-E_{gg}(q,r)\big)\big(E_g(k)-E_g(p)\big)}\,.
\end{align}
The factor $\frac12$ is the weight of the two-gluon intermediate states (Section~\ref{sec:fock}); the two terms of Eq.~\eqref{eq:m2} give equal contributions and cancel it. The first matrix element is the complex conjugate of Eq.~\eqref{eq:m3}. Since the emitted gluon is a soft mode, $r^+=p^+-k^+>\Lambda$: this contribution exists in region~1 of Eq.~\eqref{eq:ranges} only. Inserting the matrix elements (Appendix~\ref{app:onegluon}, Eq.~\eqref{eq:psi1}) we find, with the momentum fraction $\xi\equiv p^+/k^+$,
\begin{align}\label{eq:psi1_xi}
\ket{(\Psi^a_i(k))^1_{g\rho}}=\ &-\frac{ig^2f^{abc}}{4\pi^3}\int_{1+\Lambda/k^+}^{\vee/k^+}d\xi\int d^2p\ \frac{\sqrt\xi}{1-\xi}\,\frac{(p-k)^m\,(p-\xi k)^l}{(\v p-\v k)^2\,(\v p^2-\xi\v k^2)}\,\rho^c(\v k-\v p)\notag\\
&\times\Big(\delta_{il}\delta_{jm}-\frac1\xi\,\delta_{jl}\delta_{im}-\frac{1}{1-\xi}\,\delta_{ij}\delta_{lm}\Big)\ket{g^b_j(\xi k^+,\v p)}\,.
\end{align}
In the second contribution (Figure~\ref{fig:g_g_one}c) the incoming gluon splits into the gluons $p$ and $r=k-p$, and the gluon $r$ is absorbed by the valence charge. The first denominator is $E_g(k)-E_{gg}(p,r)$, which is negative definite:
\begin{align}\label{eq:psi2_def}
\ket{(\Psi^a_i(k))^2_{g\rho}}=\frac12\int d^3p\,d^3q\,d^3r\ \ket{g^b_j(p)}\,
\frac{\braket{g^b_j(p)|H_g|g^c_l(q)g^d_m(r)}\braket{g^c_l(q)g^d_m(r)|H_{ggg}|g^a_i(k)}}{\big(E_g(k)-E_{gg}(q,r)\big)\big(E_g(k)-E_g(p)\big)}\,.
\end{align}
Here $r^+=k^+-p^+>\Lambda$, i.e.\ region~2. The result is (Appendix~\ref{app:onegluon}, Eq.~\eqref{eq:psi2})
\begin{align}\label{eq:psi2_xi}
\ket{(\Psi^a_i(k))^2_{g\rho}}=\ &\frac{ig^2f^{abc}}{4\pi^3}\int_{\Lambda/k^+}^{1-\Lambda/k^+}d\xi\ \frac{\xi^{3/2}}{1-\xi}\int d^2p\ \frac{(k-p)^m\,(p-\xi k)^l}{(\v p^2-\xi\v k^2)\,(\v p-\xi\v k)^2}\,\rho^c(\v k-\v p)\notag\\
&\times\Big(\delta_{il}\delta_{jm}-\frac1\xi\,\delta_{jl}\delta_{im}-\frac{1}{1-\xi}\,\delta_{lm}\delta_{ij}\Big)\ket{g^b_j(\xi k^+,\v p)}\,.
\end{align}
The overall signs of Eqs.~\eqref{eq:psi1} and~\eqref{eq:psi2} are fixed by the matrix elements~\eqref{eq:m2} and~\eqref{eq:m3}; we have checked them numerically by evaluating Eqs.~\eqref{eq:psi1_def} and~\eqref{eq:psi2_def} directly, and through the unitarity relation~\eqref{eq:B3unit}.

Each of the four terms linear in $\rho$ has a double pole at $p^+=k^+$, which alone would give a contribution proportional to $k^+/\Lambda$. It cancels between the instantaneous and the non-instantaneous term in each of the two regions~\eqref{eq:ranges} (Appendix~\ref{app:doublepole}).

\paragraph{Two $\rho$ part.} At second order of Eq.~\eqref{eq:pt2} there are two contributions with two valence vertices (Figure~\ref{fig:g_g_two}). They differ by the order of the two vertices in light-cone time, and since the charges do not commute, also by the order of the charges. In the first one the incoming gluon is absorbed and a gluon is emitted afterwards; the intermediate state is the soft vacuum, with the denominator $E_g(k)$,
\begin{align}\label{eq:psi1rr}
\ket{(\Psi^a_i(k))^1_{g\rho\rho}}&=\int_\Lambda^\vee dp^+\!\int d^2p\ \ket{g^b_j(p)}\,\frac{\braket{g^b_j(p)|H_g|0}\braket{0|H_g|g^a_i(k)}}{E_g(k)\,\big(E_g(k)-E_g(p)\big)}\notag\\
&=\frac{g^2}{4\pi^3}\int_\Lambda^\vee dp^+\!\int d^2p\ \sqrt{\frac{k^+}{p^+}}\ \frac{k^ip^j\,\rho^b(-\v p)\,\rho^a(\v k)}{\v k^2\,(p^+\v k^2-k^+\v p^2)}\,\ket{g^b_j(p)}\,.
\end{align}
In the second one a gluon is emitted first and the incoming gluon is absorbed afterwards; the intermediate state contains the two gluons $k$ and $p$, with the denominator $E_g(k)-E_{gg}(k,p)=-E_g(p)$,
\begin{align}\label{eq:psi2rr}
\ket{(\Psi^a_i(k))^2_{g\rho\rho}}&=\frac12\int d^3p\,d^3q\,d^3r\ \ket{g^b_j(p)}\,\frac{\braket{g^b_j(p)|H_g|g^c_l(q)g^d_m(r)}\braket{g^c_l(q)g^d_m(r)|H_g|g^a_i(k)}}{\big(E_g(k)-E_{gg}(q,r)\big)\big(E_g(k)-E_g(p)\big)}\notag\\
&=-\frac{g^2}{4\pi^3}\int_\Lambda^\vee dp^+\!\int d^2p\ \sqrt{\frac{p^+}{k^+}}\ \frac{k^ip^j\,\rho^a(\v k)\,\rho^b(-\v p)}{\v p^2\,(p^+\v k^2-k^+\v p^2)}\,\ket{g^b_j(p)}\,.
\end{align}
In both expressions the later vertex stands to the left: $\rho^b(-\v p)$ of the emission to the left of $\rho^a(\v k)$ of the absorption in Eq.~\eqref{eq:psi1rr}, and the opposite order in Eq.~\eqref{eq:psi2rr}. In Eq.~\eqref{eq:psi2rr} we have dropped the disconnected term in which the emitted gluon is reabsorbed and the incoming gluon is a spectator; it is proportional to $\ket{g^a_i(k)}$ and is excluded from the sum over $\ket n\neq\ket i$ in Eq.~\eqref{eq:pt2}. Both gluons are soft modes, $\Lambda<p^+<\vee$, without a restriction on $p^+-k^+$. With the decomposition~\eqref{eq:decomp} the energy denominator cancels in the anticommutator part (Appendix~\ref{app:tworho}), and
\begin{align}\label{eq:rrsum}
\ket{(\Psi^a_i(k))^1_{g\rho\rho}}+\ket{(\Psi^a_i(k))^2_{g\rho\rho}}=\ket{(\Psi^a_i(k))_{g\{\rho,\rho\}}}+\ket{(\Psi^a_i(k))_{g[\rho,\rho]}}\,,
\end{align}
with
\begin{align}
\ket{(\Psi^a_i(k))_{g\{\rho,\rho\}}}&=-\frac{g^2}{8\pi^3}\int_\Lambda^\vee\frac{dp^+}{\sqrt{k^+p^+}}\int d^2p\ \frac{k^ip^j}{\v k^2\v p^2}\,\big\{\rho^a(\v k),\rho^b(-\v p)\big\}\ket{g^b_j(p)}\,,\label{eq:psi_anti}\\
\ket{(\Psi^a_i(k))_{g[\rho,\rho]}}&=\frac{ig^2f^{abc}}{8\pi^3}\int_\Lambda^\vee\frac{dp^+}{\sqrt{k^+p^+}}\int d^2p\ \frac{k^ip^j}{\v k^2\v p^2}\,\frac{k^+\v p^2+p^+\v k^2}{k^+\v p^2-p^+\v k^2}\,\rho^c(\v k-\v p)\ket{g^b_j(p)}\,.\label{eq:psi_comm}
\end{align}
The anticommutator part is the product of the amplitude for the absorption of the incoming gluon and the amplitude for the emission of a Weizs\"acker--Williams gluon, symmetrised in the order of the two; it is regular at $E_g(p)=E_g(k)$. The commutator part is linear in $\rho$ and has the principal-value denominator. \begin{figure}[t]
\centering
\begin{tabular}{cc}
\scalebox{0.8}{\tikzfig{gtovactog_tworho}}&\scalebox{0.8}{\tikzfig{gg_tworho}}\\[2mm]
(a)&(b)
\end{tabular}
\caption{Transition of the incoming gluon into a gluon of different momentum with two valence vertices: (a) absorption followed by emission, Eq.~\eqref{eq:psi1rr}; (b) emission followed by absorption, Eq.~\eqref{eq:psi2rr}.}
\label{fig:g_g_two}
\end{figure}

\paragraph{The coefficient $\mf B_3$.} Comparing the sum of Eqs.~\eqref{eq:psi_inst}, \eqref{eq:psi1}, \eqref{eq:psi2}, \eqref{eq:psi_anti} and~\eqref{eq:psi_comm} with the third term of Eq.~\eqref{eq:Omegag}, i.e.\ multiplying the coefficient of $\ket{g^b_j(p)}$ by $(2\pi)^3$, we obtain
\begin{align}\label{eq:B3mom}
\mf B^{ba}_{3ji}(k,p)=\ &\frac{ig^2f^{abc}\sqrt{p^+k^+}\,(k^++p^+)}{(p^+-k^+)^2}\,\frac{\rho^c(\v k-\v p)\,\delta_{ij}}{p^+\v k^2-k^+\v p^2}\,\big[\Theta_>+\Theta_<\big]
\notag\\&
+\frac{2ig^2f^{abc}\rho^c(\v k-\v p)\,(k-p)^l\sqrt{k^+p^+}}{(k^+-p^+)^2}\,\frac{k^mp^+-k^+p^m}{\v k^2p^+-\v p^2k^+}
\notag\\&\quad\times\Big[\frac{\delta_{jm}\delta_{il}}{p^+}-\frac{\delta_{im}\delta_{jl}}{k^+}-\frac{\delta_{lm}\delta_{ij}}{p^+-k^+}\Big]
\notag\\&\quad\times
\Bigg[\frac{p^+-k^+}{(\v k-\v p)^2}\,\Theta_>+\frac{\Theta_<}{\frac{\v k^2}{k^+}-\frac{(\v k-\v p)^2}{k^+-p^+}-\frac{\v p^2}{p^+}}\Bigg]
\notag\\&
-\frac{g^2k^ip^j}{\v k^2\v p^2\sqrt{k^+p^+}}\Big[\{\rho^a(\v k),\rho^b(-\v p)\}-if^{abc}\rho^c(\v k-\v p)\,\frac{k^+\v p^2+p^+\v k^2}{k^+\v p^2-p^+\v k^2}\Big]\,,
\end{align}
with
\begin{align}
\Theta_>\equiv\Theta(p^+-k^+-\Lambda)\,,\qquad \Theta_<\equiv\Theta(k^+-\Lambda-p^+)\,,
\end{align}
and $\Lambda<k^+,p^+<\vee$. The first argument of $\mf B_3$ is the momentum of the incoming gluon, with the indices $(a,i)$, and the second that of the outgoing gluon, with the indices $(b,j)$. The first line is the instantaneous interaction, the second to fourth lines are the terms with the three-gluon vertex, and the last line contains the two terms with two valence vertices. The denominators $p^+\v k^2-k^+\v p^2$ are principal values.

\subsubsection{Two outgoing gluons}\label{sec:twoout}

At order $g$ there are two contributions with two gluons in the final state (Figure~\ref{fig:g_gg}).

\paragraph{$\rho$ independent part.} The incoming gluon splits into two soft gluons through the three-gluon vertex. This is the first-order term of Eq.~\eqref{eq:pt2} with the two-gluon states as $\ket n$,
\begin{align}\label{eq:psigg_def}
\ket{(\Psi^a_i(k))_{gg}}=\frac12\int d^3p\,d^3q\ \ket{g^c_l(q)\,g^b_j(p)}\,\frac{\braket{g^c_l(q)\,g^b_j(p)|H_{ggg}|g^a_i(k)}}{E_g(k)-E_{gg}(p,q)}\,.
\end{align}
The $\delta$-function of Eq.~\eqref{eq:m3} sets $q=k-p$. With $\xi=p^+/k^+$ the three brackets of Eq.~\eqref{eq:m3} become $2(p-\xi k)^i$, $-\frac2\xi(p-\xi k)^j$ and $-\frac{2}{1-\xi}(p-\xi k)^l$, and the energy denominator is $E_g(k)-E_{gg}(p,k-p)=-(\v p-\xi\v k)^2/[2\xi(1-\xi)k^+]$, which never vanishes. Hence
\begin{align}\label{eq:psigg}
\ket{(\Psi^a_i(k))_{gg}}=\ &-\frac{igf^{abc}\sqrt{k^+}}{4\pi^{3/2}}\int_{\Lambda/k^+}^{1-\Lambda/k^+}d\xi\,\sqrt{\xi(1-\xi)}\int d^2p\ \frac{p^m-\xi k^m}{(\v p-\xi\v k)^2}\notag\\
&\times\Big[\delta_{im}\delta_{jl}-\frac1\xi\,\delta_{jm}\delta_{il}-\frac{1}{1-\xi}\,\delta_{lm}\delta_{ij}\Big]\ket{g^c_l\big((1-\xi)k^+,\v k-\v p\big)\,g^b_j(\xi k^+,\v p)}\,.
\end{align}
Both gluons are soft modes, which gives the limits of the $\xi$ integral. The integrand is symmetric under the exchange of the two gluons, $(b,j,\xi,\v p)\leftrightarrow(c,l,1-\xi,\v k-\v p)$. Comparing with the last term of Eq.~\eqref{eq:Omegag},
\begin{align}\label{eq:C1mom}
\mf C^{bca}_{1jki}(k,p)=\ &-\frac{igf^{abc}\,(p^mk^+-p^+k^m)\sqrt{k^+p^+(k^+-p^+)}}{\sqrt2\,(k^+\v p-p^+\v k)^2}\notag\\&\times\Big[\frac{\delta_{im}\delta_{jk}}{k^+}-\frac{\delta_{ik}\delta_{mj}}{p^+}-\frac{\delta_{km}\delta_{ij}}{k^+-p^+}\Big]\,\Theta(k^+-\Lambda-p^+)\,,
\end{align}
with $p^+>\Lambda$. In this coefficient $k$ is the momentum of the incoming gluon $(a,i)$, $p$ the momentum of the outgoing gluon $(b,j)$, and the second outgoing gluon $(c,k)$ carries $k-p$; in the bracket the polarisation index $k$ is not to be confused with the momentum. $\mf C_1$ is independent of $\rho$ and obeys the symmetry relation~\eqref{eq:symm}.

\paragraph{One $\rho$ part.} The valence charge emits a gluon, and the incoming gluon is a spectator. With Eq.~\eqref{eq:m2},
\begin{align}\label{eq:psiggrho}
\ket{(\Psi^a_i(k))_{gg\rho}}&=\frac12\int d^3p\,d^3q\ \ket{g^c_l(q)\,g^b_j(p)}\,\frac{\braket{g^c_l(q)\,g^b_j(p)|H_g|g^a_i(k)}}{E_g(k)-E_{gg}(p,q)}\notag\\
&=-\frac{g}{2\pi^{3/2}}\int_\Lambda^\vee\frac{dp^+}{\sqrt{p^+}}\int d^2p\ \frac{p^j}{\v p^2}\,\rho^b(-\v p)\,\ket{g^a_i(k)\,g^b_j(p)}\,,
\end{align}
where the two terms of Eq.~\eqref{eq:m2} are equal and the denominator is $E_g(k)-E_{gg}(k,p)=-E_g(p)$. This is the leading-order wave function~\eqref{eq:PsiLO} in the presence of the incoming gluon. Comparing with the fourth term of Eq.~\eqref{eq:Omegag} reproduces $\mf A^{(1)}_2$ of Eq.~\eqref{eq:A1mom}, as it must: the coefficient of $a^\dagger$ in $\Omega$ does not depend on the state on which $\Omega$ acts.

\begin{figure}[t]
\centering
\begin{tabular}{cc}
\scalebox{0.8}{\tikzfig{g_togg}}&\scalebox{0.8}{\tikzfig{gtog_plusg}}\\[2mm]
(a)&(b)
\end{tabular}
\caption{Two gluons in the final state at order $g$: (a) splitting of the incoming gluon, Eq.~\eqref{eq:psigg_def}; (b) emission of a gluon by the valence charge with the incoming gluon as a spectator, Eq.~\eqref{eq:psiggrho}.}
\label{fig:g_gg}
\end{figure}

\subsubsection{Three outgoing gluons}\label{sec:threeout}

Three gluons appear in the dressed one-gluon state at order $g^2$: through the four-gluon vertex, through two successive splittings, through a splitting accompanied by an emission from the valence charge in the two time orderings, and through two emissions from the valence charge. In the expansion~\eqref{eq:OmegaExp} they correspond to the term $\mf B_2\,a^\dagger a^\dagger$, with the incoming gluon as a spectator, and to terms with three creation operators and one annihilation operator, which are among the terms denoted by dots. These components were computed in~\cite{Radhakrishnan:2026soft}. They are not needed here, and we do not reproduce them. The reason is a counting of orders. As shown in Section~\ref{sec:observable_def}, the cross section at order $g^4$ requires the overlap of $\Omega\ket0$ after the scattering with a dressed one-gluon state to order $g^3$ and with a dressed two-gluon state to order $g^2$. The three-gluon component of $\Omega\ket0$ is of order $g^3$, and the three-gluon component of $\Omega\ket{g^a_i(k)}$ is of order $g^2$; their product is of order $g^5$. In the same way the three-gluon components of a dressed two-gluon state, of order $g$, multiply the three-gluon component of $\Omega\ket0$ and contribute to the amplitude at order $g^4$. For the same reason neither the two-gluon component of $\Omega\ket{g^a_i(k)}$ at order $g^3$ nor the four-gluon vertex $H_{gggg}$ is needed: all components of the dressed one-gluon state that enter are those computed in Sections~\ref{sec:vacout}--\ref{sec:twoout}, and $\mf A^{(3)}_1$.

\NOTE{The three-gluon components are not derived here; it is assumed, as in the calculation file, that the normal-ordered $H_{gggg}$ gives no contribution to $\mf A^{(3)}_2$. To be confirmed by the authors.}

\subsection{Summary of the wave function}\label{sec:summarywf}

Collecting the components, the dressed one-gluon state to order $g^2$, restricted to at most two gluons, is
\begin{align}\label{eq:psi_sum}
\ket{\Psi^a_i(k)}=\ &\mf N\,\ket{g^a_i(k)}+\int\frac{d^3p}{(2\pi)^3}\,\Big[\mf B^{ba}_{3ji}(k,p)\Big]_{\rho^0}\ket{g^b_j(p)}\notag\\
&+\ket{(\Psi^a_i(k))_\rho}+\ket{(\Psi^a_i(k))_{gg}}+\ket{(\Psi^a_i(k))_{gg\rho}}\notag\\
&+\ket{(\Psi^a_i(k))^1_{g\rho\,{\rm inst}}}+\ket{(\Psi^a_i(k))^2_{g\rho\,{\rm inst}}}+\ket{(\Psi^a_i(k))^1_{g\rho}}+\ket{(\Psi^a_i(k))^2_{g\rho}}\notag\\
&+\ket{(\Psi^a_i(k))_{g\{\rho,\rho\}}}+\ket{(\Psi^a_i(k))_{g[\rho,\rho]}}+\dots
\end{align}
The second line contains the components of order $g$, Eqs.~\eqref{eq:psi_rho}, \eqref{eq:psigg} and~\eqref{eq:psiggrho}; the third and fourth lines the components of order $g^2$ with one gluon, Eqs.~\eqref{eq:psi_inst}, \eqref{eq:psi1}, \eqref{eq:psi2}, \eqref{eq:psi_anti} and~\eqref{eq:psi_comm}. The first line is the component along the incoming gluon itself, the square bracket of Eq.~\eqref{eq:pt2}. It is fixed by the normalisation
\begin{align}\label{eq:normcond}
\braket{\Psi^{a'}_{i'}(k')|\Psi^a_i(k)}=\delta^{aa'}\delta_{ii'}\,\delta^{(3)}(k-k')\,,
\end{align}
which we now verify to order $g^2$. The terms of order $g^2$ in the overlap are the overlaps of two components of order $g$ and the overlaps of $\ket{g}$ with the components of order $g^2$.

The verification is carried out in Appendix~\ref{app:norm}; the result is the following. (i) The terms with two charges that are not proportional to $\delta^{(3)}(k-k')$ cancel among the overlaps of $\ket{(\Psi)_\rho}$, $\ket{(\Psi)_{gg\rho}}$ and $\ket{(\Psi)_{g\{\rho,\rho\}}}$, for non-commuting charges; the remaining term fixes the coefficient of $\ket{g^a_i(k)}$ to be $\mf N=1+\mf N_2$ with the same $\mf N_2$ as in the vacuum sector, Eq.~\eqref{eq:N2mom}. (ii) The terms linear in $\rho$ cancel; for the terms with the three-gluon vertex this requires the overlaps of the two-gluon components, Eq.~\eqref{eq:ov4}. (iii) The $\rho$-independent probability for the incoming gluon to split, Eq.~\eqref{eq:ov5}, is compensated by a term of $\mf B_3$ proportional to $\delta^{(3)}(k-p)$, Eq.~\eqref{eq:B3rho0}, the wave-function renormalisation of the soft gluon. It multiplies the scaleless integral $\int d^2\tilde p/\tilde{\v p}^2$ and vanishes in dimensional regularisation.

\NOTE{Setting the term~\eqref{eq:B3rho0} to zero relies on $\int d^2\tilde p/\tilde{\v p}^2=0$. If ultraviolet and collinear divergences are treated separately in the cross section, it has to be kept. To be decided by the authors.}

\subsection{Unitarity of \texorpdfstring{$\Omega$}{Omega} and the relations between the coefficients}\label{sec:unitarity}

The coefficients of Eq.~\eqref{eq:OmegaExp} are not independent. Since $\Omega$ maps an orthonormal basis onto an orthonormal basis,
\begin{align}\label{eq:unit}
\Omega^\dagger\,\Omega=1\,.
\end{align}
We impose this condition order by order in $g$: $\Omega^\dagger\Omega$ is brought to normal order with the commutator~\eqref{eq:aadagger}, keeping the order of the coefficients, and the coefficient of each normal-ordered product of $a^\dagger$ and $a$ is set to zero. The algebra is given in Appendix~\ref{app:unitarity}. At order $g$ and $g^2$ one finds
\begin{align}\label{eq:unit_summary}
&\mf A^{(1)\dagger b}_{1j}(p)=-\mf A^{(1)b}_{2j}(p)\,,\qquad
\mf C^{\dagger dcb}_{2lkj}(p,q)=-\mf C^{dcb}_{1lkj}(p,q)\,,\notag\\
&\mf N_2=-\frac12\int\frac{d^3p}{(2\pi)^3}\,\mf A^{(1)\dagger b}_{2j}(p)\,\mf A^{(1)b}_{2j}(p)\,,
\end{align}
together with Eq.~\eqref{eq:B1unit}, which expresses $\mf B_1$ through $\mf B_2$, $\mf C_1\mf A^{(1)}_2$ and the symmetrised product of two $\mf A^{(1)}_2$, and Eq.~\eqref{eq:B3unit}, a relation between $\mf B_3$, $\mf B_3^\dagger$ and products of $\mf A^{(1)}_2$ and $\mf C_1$. At order $g^3$ the coefficient of $a^\dagger$ gives
\begin{align}\label{eq:A31unit}
\mf A^{(3)\dagger b}_{1j}(p)=\ &-\mf A^{(3)b}_{2j}(p)-\big[\mf N_2,\mf A^{(1)b}_{2j}(p)\big]
+\int\frac{d^3q}{(2\pi)^3}\Big[2\,\mf A^{(1)c}_{1k}(q)\,\mf B^{cb}_{2kj}(p,q)\notag\\
&-\mf B^{\dagger cb}_{3kj}(p,q)\,\mf A^{(1)c}_{2k}(q)
-2\,\mf C^{\dagger dcb}_{1lkj}(p,q)\,\mf B^{cd}_{2kl}(q,p-q)\Big]\,.
\end{align}
The first two relations of Eq.~\eqref{eq:unit_summary} are satisfied by Eqs.~\eqref{eq:A1mom} and~\eqref{eq:A11mom} and determine $\mf C_2$; the third reproduces Eq.~\eqref{eq:N2mom}. Eq.~\eqref{eq:B3unit} is satisfied by Eq.~\eqref{eq:B3mom} and fixes the relative sign of its term $\propto\Theta_>$ and $\mf C_1$.

\paragraph{The order of the operators.} For commuting charges the commutator in Eq.~\eqref{eq:A31unit} is absent. Here $\mf N_2$ is bilinear in the charges at two different points and does not commute with $\mf A^{(1)}_2$; the commutator, Eq.~\eqref{eq:commN2A}, is a term with two charges $\propto g^3f^{abc}\{\rho^b,\rho^c\}\log(\vee/\Lambda)$, equal to the term $-2k^i\log(\vee/\Lambda)$ of $[\mf A^{(3)}_2]_{\rho\rho}$. The order of the charges in the three-$\rho$ term~\eqref{eq:A3threerho} is therefore not a matter of convention, and with that order
\begin{align}\label{eq:N2Acomb}
-\big[\mf A^{(3)}_2\big]_{\rho\rho\rho}-\big[\mf N_2,\mf A^{(1)}_2\big]=-\mf N_2\,\mf A^{(1)}_2\,.
\end{align}

\paragraph{What is fixed by what.} The diagrammatic calculation fixes $\mf N$, $\mf A_2=\mf A^{(1)}_2+\mf A^{(3)}_2$ and $\mf B_2$ (vacuum initial state) and $\mf A^{(1)}_1$, $\mf B_3$ and $\mf C_1$ (one-gluon initial state). Unitarity fixes $\mf C_2$, $\mf B_1$ and $\mf A^{(3)}_1$ through Eqs.~\eqref{eq:C2unit}, \eqref{eq:B1unit} and~\eqref{eq:A31unit}. The relations~\eqref{eq:A11unit}, \eqref{eq:N2unit} and~\eqref{eq:B3unit} involve only coefficients fixed by the diagrammatic calculation and are checks. We have verified the relations of this subsection in two ways: by a symbolic normal ordering of $\Omega^\dagger\Omega$ with non-commuting coefficients for a discrete set of modes, and, for the perturbative series of Section~\ref{sec:lcpt} on which they rest, by a numerical construction of $\Omega$ with matrices in place of the charges.

\subsection{Summary of the coefficients}\label{sec:coefficients}

Table~\ref{tab:coeffs} lists the coefficients of $\Omega$, the order at which they are needed, what determines them, and the equations in which they are given. After the relations~\eqref{eq:A11unit} and~\eqref{eq:C2unit} have been used, the expansion~\eqref{eq:OmegaExp} reads
\begin{align}\label{eq:OmegaFinal}
\Omega = \ &\mf{N}
+\int\frac{d^3p}{(2\pi)^3}\Big[\mf{A}^b_{j}(p)\,a^{\dagger b}_j(p)-\mf{A}^{\dagger b}_{j}(p)\,a^{b}_j(p)+\mf{A}^{(3)b}_{2j}(p)\,a^{\dagger b}_j(p)+\mf{A}^{(3)b}_{1j}(p)\,a^{b}_j(p)\Big]
\notag\\&
+\int\frac{d^3p}{(2\pi)^3}\frac{d^3q}{(2\pi)^3}\Big[\mf{B}^{cb}_{1kj}(p,q)\,a^b_j(p)a^c_k(q)+\mf{B}^{cb}_{2kj}(p,q)\,a^{\dagger b}_j(p)a^{\dagger c}_k(q)
\notag\\&\hspace{3.4cm}+\mf{B}^{cb}_{3kj}(p,q)\,a^{\dagger c}_k(q)a^b_j(p)\Big]
\notag\\&
+\int\frac{d^3p}{(2\pi)^3}\frac{d^3q}{(2\pi)^3}\Big[\mf{C}^{dcb}_{lkj}(p,q)\,a^{\dagger d}_l(q)a^{\dagger c}_k(p-q)a^b_j(p)
\notag\\&\hspace{3.4cm}-\mf{C}^{\dagger dcb}_{lkj}(p,q)\,a^{\dagger b}_j(p)a^c_k(p-q)a^d_l(q)\Big]+\dots,
\end{align}
with $\mf A\equiv\mf A^{(1)}_2$ and $\mf C\equiv\mf C_1$.

\begin{table}[t]
\centering
\renewcommand{\arraystretch}{1.3}
\footnotesize\setlength{\tabcolsep}{4pt}
\begin{tabular}{@{}l c >{\raggedright\arraybackslash}p{0.37\textwidth} c c@{}}
\toprule
coefficient & order & fixed by & momentum space & coordinate space\\
\midrule
$\mf A^{(1)}_2\equiv\mf A$ & $g$ & $\Omega\ket0$, one gluon; first order in $H_g$ & \eqref{eq:A1mom} & \eqref{eq:A1X}\\
$\mf A^{(1)}_1$ & $g$ & $\Omega\ket g$, vacuum; check: unitarity, Eq.~\eqref{eq:A11unit} & \eqref{eq:A11mom} & \eqref{eq:A1X}\\
$\mf C_1\equiv\mf C$ & $g$ & $\Omega\ket g$, two gluons; first order in $H_{ggg}$ & \eqref{eq:C1mom} & \eqref{eq:C1X}\\
$\mf C_2$ & $g$ & unitarity, Eq.~\eqref{eq:C2unit} & \eqref{eq:C2unit} & \eqref{eq:unitX}\\
$\mf N_2$ & $g^2$ & norm of $\Omega\ket0$ and of $\Omega\ket g$; unitarity, Eq.~\eqref{eq:N2unit} & \eqref{eq:N2mom} & \eqref{eq:N2X}\\
$\mf B_2$ & $g^2$ & $\Omega\ket0$, two gluons; second order~\cite{Lublinsky:2016meo} & \eqref{eq:B2mom} & \eqref{eq:B2X}\\
$\mf B_3$ & $g^2$ & $\Omega\ket g$, one gluon; second order, denominators $E_g(k)-E_n$; check: Eq.~\eqref{eq:B3unit} & \eqref{eq:B3mom} & \eqref{eq:B3X}\\
$\mf B_1$ & $g^2$ & unitarity, Eq.~\eqref{eq:B1unit} & \eqref{eq:B1unit} & \eqref{eq:unitX}\\
$\mf A^{(3)}_2$ & $g^3$ & $\Omega\ket0$, one gluon; third order~\cite{Lublinsky:2016meo} & \eqref{eq:A3onerho}--\eqref{eq:A3threerho} & \eqref{eq:A3onerhoX}--\eqref{eq:A3threerhoX}\\
$\mf A^{(3)}_1$ & $g^3$ & unitarity, Eq.~\eqref{eq:A31unit} & \eqref{eq:A31unit} & \eqref{eq:unitX}\\
\bottomrule
\end{tabular}
\caption{The coefficients of the evolution operator $\Omega$, Eq.~\eqref{eq:OmegaExp}, that are needed for the single inclusive gluon cross section at next-to-leading order. $\Omega\ket0$ and $\Omega\ket g$ denote the matching to the dressed vacuum (Section~\ref{sec:vacin}) and to the dressed one-gluon state (Section~\ref{sec:onegin}).}
\label{tab:coeffs}
\end{table}

The expansion of $\Omega$ in transverse coordinate space, Eq.~\eqref{eq:OmegaX}, and the coefficients in coordinate space, in which they are used in the cross section, are given in Appendix~\ref{app:coefficients}.

With the coefficients of Table~\ref{tab:coeffs} the operator $\Omega$ is known to the accuracy required. In Section~\ref{sec:cross} we use it to set up the calculation of the single inclusive gluon cross section.
